\subsection{Lazy Stream Evaluation: Generators}

A generator is a function-like object that produces values over time instead of computing and returning one final value immediately. The source-level difference is the presence of \texttt{yield}. Once a function body contains \texttt{yield}, calling that function does not execute the body immediately. It creates a generator object.

That generator object can later be resumed. Each resumption continues execution until the next \texttt{yield}, a \texttt{return}, or the end of the function body.

So a generator is not simply a function that returns a list slowly. It is a suspended execution mechanism.

\subsubsection{The \texttt{yield} Keyword Architecture: Pausing Function Execution Without Destroying Local State}

A normal function starts, runs, and terminates when it reaches \texttt{return} or the end of the body.

\begin{verbatim}
def normal_func():
    print("normal function starts")
    return "D'oh!"

result = normal_func()

print(f"result:       {result!r}")
print(f"type(result): {type(result)}")
\end{verbatim}

Possible output:

\begin{verbatim}
normal function starts
result:       "D'oh!"
type(result): <class 'str'>
\end{verbatim}

A generator function contains \texttt{yield}.

\begin{verbatim}
def generator_func():
    print("generator body starts")
    yield "D'oh!"

gen = generator_func()

print(f"gen:       {gen}")
print(f"type(gen): {type(gen)}")
\end{verbatim}

Possible output:

\begin{verbatim}
gen:       <generator object generator_func at 0x...>
type(gen): <class 'generator'>
\end{verbatim}

Notice that \texttt{"generator body starts"} was not printed yet. Calling \texttt{generator\_func()} created a generator object. It did not run the function body immediately.

The body starts when the generator is advanced.

\begin{verbatim}
def generator_func():
    print("generator body starts")
    yield "D'oh!"

gen = generator_func()

print("before next()")
value = next(gen)
print("after next()")
print(f"value: {value!r}")
\end{verbatim}

Possible output:

\begin{verbatim}
before next()
generator body starts
after next()
value: "D'oh!"
\end{verbatim}

The \texttt{yield} statement sends one value outward and pauses the generator frame. It does not destroy the local state.

\begin{verbatim}
def count_quotes():
    print("step 1")
    yield "D'oh!"

    print("step 2")
    yield "No soup for you!"

    print("step 3")
    yield "The answer is 42."

gen = count_quotes()

print(next(gen))
print(next(gen))
print(next(gen))
\end{verbatim}

Possible output:

\begin{verbatim}
step 1
D'oh!
step 2
No soup for you!
step 3
The answer is 42.
\end{verbatim}

Each call to \texttt{next(gen)} resumes the same suspended execution. The function body does not restart from the top each time. It continues after the previous \texttt{yield}.

A generator can keep local variables alive between yields.

\begin{verbatim}
def counter():
    n = 40

    print(f"before first yield: n={n}")
    yield n

    n = n + 1
    print(f"before second yield: n={n}")
    yield n

    n = n + 1
    print(f"before third yield: n={n}")
    yield n

gen = counter()

print(f"first:  {next(gen)}")
print(f"second: {next(gen)}")
print(f"third:  {next(gen)}")
\end{verbatim}

Possible output:

\begin{verbatim}
before first yield: n=40
first:  40
before second yield: n=41
second: 41
before third yield: n=42
third:  42
\end{verbatim}

The local name \texttt{n} survives across suspension points. The generator frame is paused, not discarded.

The important contrast is:

\begin{verbatim}
return:
    sends a value outward
    terminates the current function frame

yield:
    sends a value outward
    suspends the current generator frame
    allows later resumption
\end{verbatim}

A generator object is inspectable.

\begin{verbatim}
def counter():
    yield 42

gen = counter()

print(f"type(gen): {type(gen)}")
print()

selected_names = [
    "__iter__",
    "__next__",
    "send",
    "throw",
    "close",
    "gi_frame",
    "gi_code",
    "gi_running",
    "gi_yieldfrom",
]

for name in selected_names:
    print(f"{name:<14} {name in dir(gen)}")
\end{verbatim}

Possible output:

\begin{verbatim}
type(gen): <class 'generator'>

__iter__       True
__next__       True
send           True
throw          True
close          True
gi_frame       True
gi_code        True
gi_running     True
gi_yieldfrom   True
\end{verbatim}

Some crucial attributes and methods are:

\begin{itemize}
    \item \texttt{\_\_next\_\_}: resumes the generator until the next yielded value;
    \item \texttt{send}: resumes the generator and can send a value into it;
    \item \texttt{throw}: injects an exception into the suspended generator;
    \item \texttt{close}: requests generator shutdown;
    \item \texttt{gi\_frame}: exposes the suspended frame, while it exists;
    \item \texttt{gi\_code}: exposes the generator's code object.
\end{itemize}

Only \texttt{\_\_next\_\_}, \texttt{gi\_frame}, and \texttt{gi\_code} are needed for the current subsection. The bidirectional methods \texttt{send}, \texttt{throw}, and \texttt{close} will be studied later.

\subsubsection{Generator Objects: Suspended Frames, Instruction Pointers, and Resumable Execution State}

A generator object contains the machinery needed to resume execution later.

\begin{verbatim}
def small_gen():
    x = "D'oh!"
    yield x

gen = small_gen()

print(f"gen:          {gen}")
print(f"gen.gi_code:  {gen.gi_code}")
print(f"gen.gi_frame: {gen.gi_frame}")
\end{verbatim}

Possible output:

\begin{verbatim}
gen:          <generator object small_gen at 0x...>
gen.gi_code:  <code object small_gen at 0x..., file "...", line 1>
gen.gi_frame: <frame at 0x..., file "...", line 1, code small_gen>
\end{verbatim}

The generator has a code object and a frame object. The code object describes the compiled instructions. The frame stores the current execution state.

The code object can be inspected.

\begin{verbatim}
def small_gen():
    x = "D'oh!"
    yield x

gen = small_gen()
code_obj = gen.gi_code

print(f"type(code_obj): {type(code_obj)}")
print()

selected_names = [
    "co_name",
    "co_varnames",
    "co_consts",
    "co_names",
    "co_code",
]

for name in selected_names:
    print(f"{name:<14} {name in dir(code_obj)}")

print()
print(f"co_name:     {code_obj.co_name}")
print(f"co_varnames: {code_obj.co_varnames}")
print(f"co_consts:   {code_obj.co_consts}")
\end{verbatim}

Possible output:

\begin{verbatim}
type(code_obj): <class 'code'>

co_name        True
co_varnames    True
co_consts      True
co_names       True
co_code        True

co_name:     small_gen
co_varnames: ('x',)
co_consts:   (None, "D'oh!")
\end{verbatim}

The frame object can also be inspected.

\begin{verbatim}
def small_gen():
    x = "D'oh!"
    yield x

gen = small_gen()
frame = gen.gi_frame

print(f"type(frame): {type(frame)}")
print()

selected_names = [
    "f_code",
    "f_locals",
    "f_globals",
    "f_lineno",
    "f_lasti",
]

for name in selected_names:
    print(f"{name:<12} {name in dir(frame)}")

print()
print(f"frame.f_code.co_name: {frame.f_code.co_name}")
print(f"frame.f_locals:       {frame.f_locals}")
print(f"frame.f_lasti:        {frame.f_lasti}")
\end{verbatim}

Possible output before first resumption:

\begin{verbatim}
type(frame): <class 'frame'>

f_code       True
f_locals     True
f_globals    True
f_lineno     True
f_lasti      True

frame.f_code.co_name: small_gen
frame.f_locals:       {}
frame.f_lasti:        0
\end{verbatim}

Before the first \texttt{next()}, the generator body has not yet assigned \texttt{x}. The local namespace display is still empty.

After the first yield, the local state exists.

\begin{verbatim}
def small_gen():
    x = "D'oh!"
    yield x

gen = small_gen()

print("before next()")
print(f"locals: {gen.gi_frame.f_locals}")

value = next(gen)

print()
print("after next()")
print(f"value:  {value!r}")
print(f"locals: {gen.gi_frame.f_locals}")
print(f"lasti:  {gen.gi_frame.f_lasti}")
\end{verbatim}

Possible output:

\begin{verbatim}
before next()
locals: {}

after next()
value:  "D'oh!"
locals: {'x': "D'oh!"}
lasti:  12
\end{verbatim}

The exact value of \texttt{f\_lasti} is version-dependent. It represents the current bytecode position. The important idea is that the frame remembers where execution stopped.

A generator can be disassembled.

\begin{verbatim}
import dis

def small_gen():
    x = "D'oh!"
    yield x
    y = "No soup for you!"
    yield y

dis.dis(small_gen)
\end{verbatim}

Possible output excerpt:

\begin{verbatim}
... RETURN_GENERATOR
... RESUME
... LOAD_CONST               ... ("D'oh!")
... STORE_FAST               ... (x)
... LOAD_FAST                ... (x)
... YIELD_VALUE
... RESUME
... LOAD_CONST               ... ('No soup for you!')
... STORE_FAST               ... (y)
... LOAD_FAST                ... (y)
... YIELD_VALUE
...
\end{verbatim}

The exact instruction list differs across Python versions, but \texttt{YIELD\_VALUE} is the critical operation. It is the point where execution produces a value and suspends.

A completed generator has no active frame.

\begin{verbatim}
def one_value():
    yield 42

gen = one_value()

print(f"before first next: {gen.gi_frame}")

print(next(gen))

try:
    next(gen)
except StopIteration:
    pass

print(f"after completion:  {gen.gi_frame}")
\end{verbatim}

Possible output:

\begin{verbatim}
before first next: <frame at 0x..., file "...", line 1, code one_value>
42
after completion:  None
\end{verbatim}

After the generator is exhausted, \texttt{gi\_frame} becomes \texttt{None}. There is no suspended execution state left to resume.

The generator object may still exist, but it is exhausted.

\begin{verbatim}
def one_value():
    yield 42

gen = one_value()

print(next(gen))

for attempt in [1, 2]:
    try:
        print(next(gen))
    except StopIteration as err:
        print(f"attempt {attempt}: {type(err).__name__}")
\end{verbatim}

Possible output:

\begin{verbatim}
42
attempt 1: StopIteration
attempt 2: StopIteration
\end{verbatim}

Once exhausted, repeated advancement raises \texttt{StopIteration} again.

\subsubsection{Execution State Resumption: Re-Entering Suspended Function Frames Across Iteration Steps}

A generator can be advanced manually with \texttt{next()} or automatically by a \texttt{for} loop.

Manual advancement shows the resumption points clearly.

\begin{verbatim}
def quotes():
    print("A")
    yield "D'oh!"

    print("B")
    yield "No soup for you!"

    print("C")
    yield "The answer is 42."

gen = quotes()

print("created generator")
print()

print(f"next 1 -> {next(gen)!r}")
print()
print(f"next 2 -> {next(gen)!r}")
print()
print(f"next 3 -> {next(gen)!r}")
\end{verbatim}

Possible output:

\begin{verbatim}
created generator

A
next 1 -> "D'oh!"

B
next 2 -> 'No soup for you!'

C
next 3 -> 'The answer is 42.'
\end{verbatim}

The generator resumes after the previous \texttt{yield}. It does not restart from the first line.

A \texttt{for} loop performs repeated \texttt{next()} calls and stops when \texttt{StopIteration} is raised.

\begin{verbatim}
def quotes():
    yield "D'oh!"
    yield "No soup for you!"
    yield "The answer is 42."

for quote in quotes():
    print(f"quote: {quote}")
\end{verbatim}

Possible output:

\begin{verbatim}
quote: D'oh!
quote: No soup for you!
quote: The answer is 42.
\end{verbatim}

The loop is roughly equivalent to this manual structure:

\begin{verbatim}
gen = quotes()

while True:
    try:
        quote = next(gen)
    except StopIteration:
        break

    print(f"quote: {quote}")
\end{verbatim}

A generator is its own iterator.

\begin{verbatim}
def quotes():
    yield "D'oh!"

gen = quotes()
iterator = iter(gen)

print(f"gen:              {gen}")
print(f"iterator:         {iterator}")
print(f"gen is iterator:  {gen is iterator}")
print(f"type(iterator):   {type(iterator)}")
\end{verbatim}

Possible output:

\begin{verbatim}
gen:              <generator object quotes at 0x...>
iterator:         <generator object quotes at 0x...>
gen is iterator:  True
type(iterator):   <class 'generator'>
\end{verbatim}

The generator object implements both \texttt{\_\_iter\_\_} and \texttt{\_\_next\_\_}. This is why it can be used directly in a \texttt{for} loop.

The generator remembers local state across loop steps.

\begin{verbatim}
def running_total(values):
    total = 0

    for value in values:
        total = total + value
        print(f"inside generator: value={value}, total={total}")
        yield total

nums = [10, 20, 12]

for total in running_total(nums):
    print(f"outside loop:     total={total}")
\end{verbatim}

Possible output:

\begin{verbatim}
inside generator: value=10, total=10
outside loop:     total=10
inside generator: value=20, total=30
outside loop:     total=30
inside generator: value=12, total=42
outside loop:     total=42
\end{verbatim}

The local variable \texttt{total} survives across yields. Each resumed step continues with the previous value of \texttt{total}.

A generator can produce an unbounded stream because it does not need to build a complete list before returning.

\begin{verbatim}
def count_up(start):
    n = start

    while True:
        yield n
        n = n + 1

gen = count_up(40)

print(next(gen))
print(next(gen))
print(next(gen))
\end{verbatim}

Possible output:

\begin{verbatim}
40
41
42
\end{verbatim}

This generator has no natural end. The caller decides how many values to request.

A list-building version would need to decide the full size in advance.

\begin{verbatim}
def make_three_numbers(start):
    numbers = []

    numbers.append(start)
    numbers.append(start + 1)
    numbers.append(start + 2)

    return numbers

def generate_numbers(start):
    yield start
    yield start + 1
    yield start + 2

list_result = make_three_numbers(40)
gen_result = generate_numbers(40)

print(f"list_result:       {list_result}")
print(f"type(list_result): {type(list_result)}")
print()
print(f"gen_result:        {gen_result}")
print(f"type(gen_result):  {type(gen_result)}")
\end{verbatim}

Possible output:

\begin{verbatim}
list_result:       [40, 41, 42]
type(list_result): <class 'list'>

gen_result:        <generator object generate_numbers at 0x...>
type(gen_result):  <class 'generator'>
\end{verbatim}

The list contains all values immediately. The generator contains suspended execution that can produce values later.

This is why generators are useful for lazy stream evaluation:

\begin{quote}
A generator produces the next value when asked, while preserving enough execution state to continue later.
\end{quote}

\subsubsection{Returning from Generators: \texttt{StopIteration} and Generator Completion Values}

A generator ends when execution reaches the end of its body or when it executes \texttt{return}.

When a generator ends, advancing it raises \texttt{StopIteration}.

\begin{verbatim}
def one_quote():
    yield "D'oh!"

gen = one_quote()

print(next(gen))

try:
    next(gen)
except StopIteration as err:
    print(f"error type: {type(err)}")
    print(f"message:    {err}")
    print(f"err.args:   {err.args}")
\end{verbatim}

Possible output:

\begin{verbatim}
D'oh!
error type: <class 'StopIteration'>
message:    
err.args:   ()
\end{verbatim}

The \texttt{StopIteration} object is the signal that the stream has no next value.

It can be inspected.

\begin{verbatim}
err = StopIteration("finished")

print(f"err:       {err}")
print(f"type(err): {type(err)}")
print()

selected_names = [
    "args",
    "value",
    "__class__",
    "__str__",
    "__repr__",
    "with_traceback",
]

for name in selected_names:
    print(f"{name:<18} {name in dir(err)}")

print()
print(f"err.args:  {err.args}")
print(f"err.value: {err.value!r}")
\end{verbatim}

Possible output:

\begin{verbatim}
err:       finished
type(err): <class 'StopIteration'>

args               True
value              True
__class__          True
__str__            True
__repr__           True
with_traceback     True

err.args:  ('finished',)
err.value: 'finished'
\end{verbatim}

A generator can use \texttt{return} with a value. That value does not appear as a normal yielded item. It becomes the \texttt{value} attribute of the \texttt{StopIteration} exception.

\begin{verbatim}
def gen_with_return():
    yield "D'oh!"
    return "finished with 42"

gen = gen_with_return()

print(next(gen))

try:
    next(gen)
except StopIteration as err:
    print(f"StopIteration.value: {err.value!r}")
    print(f"StopIteration.args:  {err.args}")
\end{verbatim}

Possible output:

\begin{verbatim}
D'oh!
StopIteration.value: 'finished with 42'
StopIteration.args:  ('finished with 42',)
\end{verbatim}

The return value is a completion value, not a stream value.

A \texttt{for} loop consumes \texttt{StopIteration} internally and does not expose the completion value.

\begin{verbatim}
def gen_with_return():
    yield "D'oh!"
    yield "No soup for you!"
    return "finished with 42"

for item in gen_with_return():
    print(f"item: {item}")
\end{verbatim}

Possible output:

\begin{verbatim}
item: D'oh!
item: No soup for you!
\end{verbatim}

The string \texttt{"finished with 42"} is not printed because it was not yielded. It was carried by \texttt{StopIteration} when the generator completed.

Manual iteration can observe it.

\begin{verbatim}
def gen_with_return():
    yield "D'oh!"
    yield "No soup for you!"
    return "finished with 42"

gen = gen_with_return()

while True:
    try:
        item = next(gen)
        print(f"yielded: {item!r}")
    except StopIteration as err:
        print(f"completion value: {err.value!r}")
        break
\end{verbatim}

Possible output:

\begin{verbatim}
yielded: "D'oh!"
yielded: 'No soup for you!'
completion value: 'finished with 42'
\end{verbatim}

A bare \texttt{return} in a generator is equivalent to completing with \texttt{None}.

\begin{verbatim}
def gen_bare_return():
    yield "Ni!"
    return

gen = gen_bare_return()

print(next(gen))

try:
    next(gen)
except StopIteration as err:
    print(f"value: {err.value!r}")
\end{verbatim}

Possible output:

\begin{verbatim}
Ni!
value: None
\end{verbatim}

Reaching the end of the generator body also completes with \texttt{None}.

\begin{verbatim}
def gen_end():
    yield "It is only a model."

gen = gen_end()

print(next(gen))

try:
    next(gen)
except StopIteration as err:
    print(f"value: {err.value!r}")
\end{verbatim}

Possible output:

\begin{verbatim}
It is only a model.
value: None
\end{verbatim}

The completion model is:

\begin{verbatim}
yield value:
    produce value to caller
    suspend generator

return final_value:
    terminate generator
    raise StopIteration(final_value) to the advancing mechanism

fall off end:
    terminate generator
    raise StopIteration(None)
\end{verbatim}

A final example shows the full lifecycle.

\begin{verbatim}
def lifecycle():
    print("start")
    yield 40

    print("middle")
    yield 41

    print("end")
    return 42

gen = lifecycle()

print(f"created: {gen}")
print()

while True:
    try:
        value = next(gen)
        print(f"caller received yield: {value}")
        print(f"frame still exists:    {gen.gi_frame is not None}")
        print()
    except StopIteration as err:
        print(f"generator completed")
        print(f"completion value: {err.value}")
        print(f"frame still exists: {gen.gi_frame is not None}")
        break
\end{verbatim}

Possible output:

\begin{verbatim}
created: <generator object lifecycle at 0x...>

start
caller received yield: 40
frame still exists:    True

middle
caller received yield: 41
frame still exists:    True

end
generator completed
completion value: 42
frame still exists: False
\end{verbatim}

The generator frame exists while the generator is suspended. After completion, the frame is gone.

The final rule is:

\begin{quote}
A generator call creates a generator object. Each \texttt{yield} pauses that object's frame and returns one stream value. A generator \texttt{return} terminates the frame and exposes its value only through \texttt{StopIteration.value}.
\end{quote}